% TODO make hw stylesheet
\documentclass[10pt]{article}
\usepackage[letterpaper,top=1in,left=1in,right=1in]{geometry}
\usepackage[dvipsnames,svgnames]{xcolor}
\usepackage[shortlabels]{enumitem}
\usepackage[framemethod=TikZ]{mdframed}
\usepackage{amsmath,amssymb,amsthm}
\usepackage[colorlinks]{hyperref}
\usepackage{multicol}
\usepackage{tabularx}

\hypersetup{
    colorlinks=true,
    linkcolor=blue,
    filecolor=magenta,
    urlcolor=magenta,
    pdftitle={MAT 324 HW}
}

\usepackage{mathtools}
\usepackage{thmtools}
\usepackage[textsize=footnotesize]{todonotes}
\usepackage{titling}
\usepackage{cancel}

\reversemarginpar
\mdfdefinestyle{mdbluebox}{roundcorner=10pt,innerbottommargin=9pt,
linecolor=blue,backgroundcolor=TealBlue!5,}
\declaretheoremstyle[headfont=\sffamily\bfseries\color{MidnightBlue},
mdframed={style=mdbluebox},]{thmbluebox}

\mdfdefinestyle{mdbloobox}{frametitlefont=\bfseries,innerbottommargin=8pt,
nobreak=true,backgroundcolor=TealBlue!5,linecolor=blue,}
\declaretheoremstyle[headfont=\bfseries\color{MidnightBlue},
mdframed={style=mdbloobox},headpunct={\\[3pt]},postheadspace=0pt,]{thmbloobox}

\mdfdefinestyle{mdredbox}{frametitlefont=\bfseries,innerbottommargin=8pt,
nobreak=true,backgroundcolor=Salmon!5,linecolor=RawSienna,}
\declaretheoremstyle[headfont=\bfseries\color{RawSienna},
mdframed={style=mdredbox},headpunct={\\[3pt]},postheadspace=0pt,]{thmredbox}

\mdfdefinestyle{mdgreenbox}{linecolor=ForestGreen,backgroundcolor=ForestGreen!5,
linewidth=2pt,rightline=false,leftline=true,topline=false,bottomline=false,}
\declaretheoremstyle[headfont=\bfseries\sffamily\color{ForestGreen!70!black},
mdframed={style=mdgreenbox},headpunct={ --- },]{thmgreenbox}

\mdfdefinestyle{mdorangebox}{linecolor=RedOrange,backgroundcolor=RedOrange!5,
linewidth=2pt,rightline=false,leftline=true,topline=false,bottomline=false,}
\declaretheoremstyle[headfont=\bfseries\sffamily\color{RedOrange!70!black},
mdframed={style=mdorangebox},headpunct={ --- },]{thmorangebox}

\mdfdefinestyle{mdblackbox}{linecolor=black,backgroundcolor=RedViolet!5!gray!5,
linewidth=3pt,rightline=false,leftline=true,topline=false,bottomline=false,}
\declaretheoremstyle[mdframed={style=mdblackbox}]{thmblackbox}

\declaretheoremstyle[spaceabove=6pt,spacebelow=6pt]{basehead}

\declaretheorem[style=basehead,name=Problem]{problem}
\declaretheorem[style=thmredbox,name=Problem,numbered=no]{problem*}
\declaretheorem[style=thmbloobox,name=Setup]{setup} % for config geo
\declaretheorem[style=thmredbox,name=Example,sibling=problem]{example}
\declaretheorem[style=thmredbox,name=$\bigstar$ Problem,sibling=problem]{niceprob}
\declaretheorem[style=thmbluebox,name=Theorem,sibling=problem]{theorem}
\declaretheorem[style=thmbluebox,name=Corollary,sibling=theorem]{corollary}
\declaretheorem[style=thmbluebox,name=Proposition,sibling=theorem]{prop}
\declaretheorem[style=thmbluebox,name=Lemma,numberwithin=problem]{lemma}
\declaretheorem[style=thmbluebox,name=Theorem,numbered=no]{theorem*}
\declaretheorem[style=thmbluebox,name=Lemma,numbered=no]{lemma*}
\declaretheorem[style=thmgreenbox,name=Claim,numberwithin=problem]{claim}
\declaretheorem[style=thmgreenbox,name=Claim,numbered=no]{claim*}
\declaretheorem[style=thmblackbox,name=Remark,numberwithin=problem]{remark}
\declaretheorem[style=thmblackbox,name=Remark,numbered=no]{remark*}
\declaretheorem[style=thmgreenbox,name=Definition,sibling=theorem]{definition}
\declaretheorem[style=thmgreenbox,name=Definition,numbered=no]{definition*}

\providecommand{\ol}{\overline}
\providecommand{\ceil}[1]{\left\lceil #1 \right\rceil}
\providecommand{\floor}[1]{\left\lfloor #1 \right\rfloor}
\newcommand{\dd}{\,\mathrm{d}}
\providecommand{\ul}{\underline}
\providecommand{\eps}{\varepsilon}
\providecommand{\half}{\frac{1}{2}}
\providecommand{\CC}{\mathbb C}
\providecommand{\FF}{\mathbb F}
\providecommand{\NN}{\mathbb N}
\providecommand{\QQ}{\mathbb Q}
\providecommand{\RR}{\mathbb R}
\providecommand{\ZZ}{\mathbb Z}
\providecommand{\dg}{^\circ}
\providecommand{\ii}{\item}
\providecommand{\setdiff}{\smallsetminus}
\providecommand{\alert}[1]{{\sffamily\textbf{\textcolor{blue}{#1}}}}
\providecommand{\inv}{^{-1}}
\providecommand{\mb}[1]{\mathbf{#1}}

\newenvironment{soln}{\begin{proof}[Solution]}{\end{proof}}

\providecommand{\pow}{\operatorname{Pow}}
\providecommand{\vocab}[1]{{\textbf{\textcolor{ForestGreen}{#1}}}}
\providecommand{\lcm}{\operatorname{lcm}}
\providecommand{\abs}[1]{\left|#1\right|}

\usepackage{mathrsfs}

% place title at top right
\setlength{\droptitle}{-8ex}
\pretitle{\begin{flushright}\Large\bfseries}
\posttitle{\par\end{flushright}}
\preauthor{\begin{flushright}}
\postauthor{\end{flushright}}
\predate{\begin{flushright}}
\postdate{\end{flushright}}

\title{MAT 324 HW05}
\author{\large Aditya Pahuja}
\date{\large Due 10/02}
\begin{document}
\maketitle

\begin{problem}
    Let $(X,\mathcal S,\mu)$ be a measure space and
    $f\colon X\to[-\infty,\infty]$ be $\mu$-integrable.
    Show that $f\ge 0$ almost everywhere on $X$ if and only if
    $\int_A f\dd\mu\ge 0$ for all $A\in\mathcal S$.
\end{problem}

\begin{soln}
    If $f\ge 0$ almost everywhere on $X$, let $A\in\mathcal S$
    and let $X' = \{x\in X : f(x) \ge 0\}$, so that $\mu(A\setminus X') = 0$
    and $f\ge 0$ on $A\cap X'$. Then
    \begin{equation*}
	\int_Af\dd\mu = \underbrace{\int_{A\cap X'}f\dd\mu}_{\ge0}
	+ \underbrace{\int_{A\setminus X'}f\dd\mu}_{=0} \ge 0.
    \end{equation*}

    Conversely suppose that $f$ is negative on
    a set $A\in\mathcal S$ of positive measure.
    By Axler's exercise 3A \#3, $\int_A-f\dd\mu > 0 \iff \int_Af < 0$,
    which proves the \emph{only if} direction.
\end{soln}

\begin{problem}
    [Axler 3B \#3]
    Suppose $\lambda$ is the Lebesgue measure on $\RR$
    and $f\colon\RR\to\RR$ is a Borel measurable function such that
    $\int\abs f\dd\lambda < \infty$. Define $g\colon\RR\to\RR$ by
    $g(x) = \int_{(-\infty,x)}f\dd\lambda$.
    Prove that $g$ is uniformly continuous on $\RR$.
\end{problem}

\begin{soln}
    The integral bound tells us $\int_X f_\pm \le \int_Xf < \infty$,
    so $\int_Xf$ is well-defined and finite.
    Let $\eps>0$. By Axler's 3.28 there exists $\delta>0$ such that
    \begin{equation*}
	g(x) - g(y) = \int_{(x,y)} f\dd\lambda < \eps
    \end{equation*}
    for all $x,y\in\RR$ such that $x\le y$ and $\lambda((x,y)) = y-x<\delta$,
    so $g$ is uniformly continuous.
\end{soln}

\pagebreak

\begin{problem}
    [Axler 3B \#4]
    \begin{enumerate}[(a)]
	\ii Suppose $(X,\mathcal S,\mu)$ is a measure space with $\mu(X) < \infty$.
	Suppose that $f\colon X\to[0,\infty)$ is
	a bounded $\mathcal S$-measurable function. Prove that
	\begin{equation*}
	    \int f\dd\mu = \inf\left\{ \sum_{j=1}^m \mu(A_j)\sup_{A_j}f :
	    A_1,\dots,A_m\text{ is an $\mathcal S$-partition of $X$}\right\}.
	\end{equation*}
	\ii Show that the conclusion of (a) can fail if
	the hypothesis that $f$ is bounded is replaced by
	the hypothesis that $\int f\dd\mu < \infty$.
	\ii Show that the conclusion of (a) can fail if
	the condition that $\mu(X) < \infty$ is deleted.
\end{enumerate}
\end{problem}

\begin{soln}
    We refer to the infimum of the upper Lebesgue sums
    as the \vocab{upper Lebesgue integral}, in contrast with
    the standard \vocab{(lower) Lebesgue integral} that we defined with lower sums.
    \\

    (a) Since $f$ is bounded, there exists $C > 0$
    such that $f(x) < C$ for all $x$. Splitting $f$ as $C - (C-f)$ yields
    \begin{align*}
	\int_Xf
	= C\mu(X) - \int_X(C-f)
	&= C\mu(X) - \sup\left\{ \sum_{j=1}^m \mu(A_j)\inf_{A_j}(C-f) :
	A_1,\dots,A_m\text{ is an $\mathcal S$-partition of $X$}\right\}\\
	&= C\mu(X) - \sup\left\{ C\mu(X) - \sum_{j=1}^m \mu(A_j)\sup_{A_j}f :
	A_1,\dots,A_m\text{ is an $\mathcal S$-partition of $X$}\right\}\\
	&= \cancel{C\mu(X) - C\mu(X)}
	+ \inf\left\{ \sum_{j=1}^m \mu(A_j)\sup_{A_j}f :
	A_1,\dots,A_m\text{ is an $\mathcal S$-partition of $X$}\right\}.
    \end{align*}

    (b) Consider $(0,1]$ with the Lebesgue measure $\lambda$
    and let $f\colon (0,1]\to\RR$ be the Lebesgue measurable function
    $f(x) = 1/\sqrt x$. Let $P$ be a Lebesgue partition of $(0,1]$
    (in analogy with ``$\mathcal S$-partition'') and
    let $P' = \{I \in P : 0\in\ol I\}$;
    $P'$ is nonempty because $0\in \ol{\bigcup_{I\in P}I} = \bigcup_{I\in P} \ol I$,
    with the equality because $P$ is finite.
    The union $A = \bigcup_{I\in P'} I$ has $\lambda(A) > 0$
    because either $A = (0,1]$ or we can find
    $I\in P\setminus P'$ with $\inf I\ne0$ minimal, so $(0,\inf I)\subseteq A$;
    thus some $J\in P'$ has $\lambda(J) > 0$.
    Since $\lim_{x\to0}f(x) = \infty$, $\sup_Jf=\infty$, so
    \begin{equation*}
	\sum_{I\in P} \lambda(I)\sup_{I}f
	\ge \lambda(J)\sup_{J}f = \infty
	\implies \inf\left\{
	    \sum_{I\in P} \lambda(I)\sup_{I}f : P\text{
		is a Lebesgue partition of $X$
	    }
	\right\} = \infty.
    \end{equation*}

    On the other hand, $f$ is the limit of the increasing sequence of functions
    $f_n = f\cdot\mathbf 1_{[1/n,1]}$ that are Riemann integrable on $[0,1]$,
    so by monotone convergence and Axler's 3.34,
    \begin{equation*}
	\int_{(0,1]} f\dd\lambda =
	\lim_{n\to\infty}\left(\int_{[1/n,1]} f_n\dd\lambda\right) =
	\lim_{n\to\infty}\left(\int_{1/n}^1 \frac1{\sqrt x}\dd x\right) =
	\lim_{n\to\infty} \left(\frac 23 -
	\frac 23\cdot\left(\frac 1n\right)^{3/2}\right) = \frac 23\ne\infty.
    \end{equation*}

    (c) Consider $[1,\infty)$ endowed with the Lebesgue measure $\lambda$
    and let $g\colon[1,\infty)\to[0,\infty)$ be the bounded Lebesgue measurable
    function $g(x) = 1/x^2$.
    Let $P$ be a Lebesgue partition of $[0,\infty)$.
    Since the union of the sets in $P$
    has infinite Lebesgue measure, some $J\in P$ has $\lambda(J) = \infty$.
    Also, $\sup_Jg > 0$ because $g> 0$, so
    \begin{equation*}
	\sum_{I\in P} \lambda(I)\sum_{I}g\ge\lambda(J)\sup_Jg = \infty
	\implies \inf\left\{
	    \sum_{I\in P} \lambda(I)\sup_{I}g : P\text{
		is a Lebesgue partition of $X$
	    }
	\right\} = \infty.
    \end{equation*}

    On the other hand, $g$ is the limit of the increasing sequence of functions
    $g_n = g\cdot\mathbf 1_{[1,n]}$ that are Riemann integrable on $[1,n]$,
    so by monotone convergence and Axler's 3.34,
    \begin{equation*}
	\int_{[1,\infty)}g\dd\lambda =
	\lim_{n\to\infty}\left(\int_{[1,n]}g_n\dd\lambda\right) =
	\lim_{n\to\infty}\left(\int_1^n\frac{1}{x^2}\dd x\right) =
	\lim_{n\to\infty}(1 - 1/n) = 1\ne\infty.\qedhere
    \end{equation*}
\end{soln}

\pagebreak

\begin{problem}
    Let $(X,\mathcal S,\mu)$ be a measure space and
    $g_1,g_2,\dots\colon X\to[0,\infty]$ be a sequence
    of $\mathcal S$-measurable functions.
    \begin{enumerate}[(a)]
	\ii Show that $\int_X(\liminf_{n\to\infty}g_n)\dd\mu
	\le \liminf_{n\to\infty}(\int_Xg_n\dd\mu)$.
	\ii Suppose also that the sequence $g_1,g_2,\dots$ converges (pointwise)
	to a function $g\colon X\to[0,\infty]$ and
	\begin{equation*}
	    \lim_{n\to\infty}\left(\int_X g_n\dd\mu\right)
	    = \int_X g\dd\mu < \infty.
	\end{equation*}
	Show that $\lim_{n\to\infty}(\int_Eg_n\dd\mu) = \int_Eg\dd\mu$
	for all $E\in\mathcal S$.
	\ii Suppose in addition (to all of the above) that
	$f_1,f_2,\dots\colon X\to[-\infty,\infty]$ is a sequence
	of $\mathcal S$-measurable functions converging to
	a function $f\colon X\to[-\infty,\infty]$ such that
	$\abs{f_n} \le g_n$ for all $n\in\NN$. Show that
	\begin{equation*}
	    \int_Xf\dd\mu = \lim_{n\to\infty}\left(\int_Xf_n\dd\mu\right).
	\end{equation*}
\end{enumerate}
\end{problem}

\begin{soln}
    (a) For any $N\in\NN$ we have that
    \begin{equation*}
	\int_X(\inf_{n\ge N}g_n)\dd\mu \le \int_X g_n\dd\mu
	\text{ for all $n\ge N$}
	\implies \int_X(\inf_{n\ge N}g_n)\dd\mu \le \inf_{n\ge N} \int_X g_n\dd\mu
    \end{equation*}
    Since $\inf_{n\ge N}g_n$ is increasing in $N$,
    we can take supremums and use monotone convergence to get
    \begin{equation*}
	\int_X \liminf_{n\to\infty} g_n
	= \sup_{N\to\infty}\left(\int_X(\inf_{n\ge N}g_n)\dd\mu\right)
	\le \sup_{N\to\infty}\inf_{n\ge N}\left(\int_Xg_n\dd\mu\right)
	= \liminf_{n\to\infty}\left(\int_Xg_n\dd\mu\right).
    \end{equation*}

    (b) It suffices to show that
    $\liminf_{n\to\infty}(\int_Eg_n\dd\mu) =
    \limsup_{n\to\infty}(\int_Eg_n\dd\mu) = \int_Eg\dd\mu$;
    to do this, we can use Fatou's lemma on the sequences
    $\int_{X\setminus E}g_n\dd\mu$ and $\int_Eg_n\dd\mu$ in the following form:
    \begin{align*}
	\liminf_{n\to\infty}\left(\int_Eg_n\dd\mu\right) \le
	\limsup_{n\to\infty}\left(\int_Eg_n\dd\mu\right)
	&= \limsup_{n\to\infty}\left(\int_Xg\dd\mu -
	\int_{X\setminus E}g_n\dd\mu\right)\\
	&= \int_Xg\dd\mu -
	\liminf_{n\to\infty}\left(\int_{X\setminus E}g_n\dd\mu\right)\\
	&\le \int_Xg\dd\mu - \int_{X\setminus E}g\dd\mu = \int_E g\dd\mu
	\le \liminf_{n\to\infty}\left(\int_Eg\dd\mu\right).
    \end{align*}

    (c) Analogous to Axler's 3.32, for each $E\in\mathcal S$ and $n\in\NN$
    we have the bound
    \begin{align*}
	\abs{\int_Xf_n - \int_Xf} &\le
	\abs{\int_{X\setminus E} f_n} + \abs{\int_{X\setminus E} f} +
	\abs{\int_E f_n - \int_E f} \le
	\int_{X\setminus E} g_n + \int_{X\setminus E}g
	+ \abs{\int_E f_n - \int_E f}.
	\tag{$\ast$}\label{eq:dct-bound}
    \end{align*}
    (for conciseness, we drop $\dd\mu$ from all the integrals here).
    \vspace{6pt}

    \textbf{Case 1.}
    Suppose that $\mu(X) < \infty$ and let $\eps>0$;
    then, let $\delta>0$ be as in Axler's proof of 3.31,
    but with the bound on $\int_Bg$ replaced by $\eps/6$.
    By (b), if $n$ is sufficiently large,
    $\abs{\int_B g_n - \int_Bg} < \eps/6$, which implies $\int_B g_n < \eps/3$
    for all $B\in\mathcal S$ with $\mu(B) < \delta$.
    Now let $E\in\mathcal S$ as in Axler's proof and use $(\ast)$ to get
    \begin{equation*}
	\tag{$\dagger$}
	\abs{\int_Xf_n - \int_Xf} \le \frac\eps2 + \abs{\int_E (f_n-f)}
    \end{equation*}
    from Axler's proof if $n$ is large enough;
    thus $\int_Xf_n$ converges to $\int_Xf$
    with the same logic as Axler's proof.
    \vspace{6pt}

    \textbf{Case 2.} If $\mu(X) = \infty$,
    pick $E\in\mathcal S$ as in Axler's proof, but with $\eps/6$ instead of $\eps/4$.
    Again by (b), if $n$ is sufficiently large, then
    $\int_{X\setminus E}g_n\dd\mu < \eps/3$,
    so applying $(\ast)$ gives the bound $(\dagger)$,
    which again implies convergence of the integrals.
\end{soln}

\pagebreak

\begin{problem}
    Let $\lambda$ denote the Lebesgue measure on $[0,\infty)$.
    For each $n\in\NN$, define
    \begin{equation*}
	f_n,g_n\colon[0,\infty)\to\RR,\quad
	f_n(x) = \frac{n^2xe^{-n^2x^2}}{1+x},\quad
	g_n(x) = \frac{xe^{-x^2}}{1 + x/n}.
    \end{equation*}
    The integrals in (a) and (b) below are improper Riemann integrals.
    \begin{enumerate}[(a)]
	\ii Find $\int_0^\infty (\lim_{n\to\infty}f_n)\dd x$ and
	$\int_0^\infty (\lim_{n\to\infty}g_n)\dd x$.
	\ii Show that
	\begin{equation*}
	    \lim_{n\to\infty}\left(\int_0^\infty f_n\dd x\right)
	    = \lim_{n\to\infty}\left(\int_0^\infty g_n\dd x\right)
	\end{equation*}
	and find this limit.
	\ii Show that there exists no Lebesgue measurable function
	$F\colon[0,\infty)\to[0,\infty]$ such that
	\begin{equation*}
	    f_n\le F\text{ a.e. for all $n\in\NN$ and }
	    \int_{[0,\infty)}F\dd\lambda<\infty.
	\end{equation*}
\end{enumerate}
\end{problem}

\begin{soln}
    (a) We can compute the limits (using l'H\^opital's rule for $f$)
    \begin{align*}
	f\coloneq\lim_{n\to\infty}f_n
	&= \frac1{1+x}\cdot\lim_{n\to\infty}\frac{n^2x}{e^{n^2x^2}}
	= \frac1{1+x}\cdot\lim_{n\to\infty}
	\frac{\cancel{n^2}}{\cancel{n^2}\cdot2x e^{n^2x^2}}
	= 0,\\
	g\coloneq\lim_{n\to\infty}g_n
	&= \frac{xe^{-x^2}}{1+\lim_{n\to\infty} x/n}
	= xe^{-x^2}.
    \end{align*}
    Then, the integrals of $f$ and $g$ are
    \begin{equation*}
	\int_0^\infty f\dd x = \boxed0,\quad
	\int_0^\infty g\dd x = \lim_{c\to\infty}\left(\int_0^c g\dd x\right)
	= \lim_{c\to\infty} \left(-\half e^{-c^2} + \half e^0\right) = \boxed\half.
    \end{equation*}

    (b) For $n\in\NN$, let $u = nx$. Then
    \begin{equation*}
	\int_0^\infty g_n\dd x
	= \int_0^\infty \frac{ue^{-u^2}}{1 + u/n}\dd u
	= \int_0^\infty \frac{nx e^{-n^2x^2}}{1 + x}\cdot n\dd x
	= \int_0^\infty f_n\dd x,
    \end{equation*}
    so $\int_0^\infty f_n\dd x$ and $\int_0^\infty g_n\dd x$ have the same limit.

    For any Lebesgue measurable function $h\colon[0,\infty)\to\RR$
    such that the Lebesgue integral $\int_{[0,\infty)}h\dd\lambda$
    and improper Riemann integral $\int_0^\infty h\dd x$
    are well-defined, the two integrals must be equal;
    to show this, we can pass the Riemann integral to the Lebesgue integral
    by Axler's 3.34 and use monotone convergence in the following fashion:
    \begin{equation*}
	\int_0^\infty h\dd x = \lim_{c\to\infty}\left(\int_0^c h\dd x\right)
	= \lim_{c\to\infty} \left(\int_{[0,c]} h\dd\lambda\right)
	= \lim_{c\to\infty} \left(\int_{[0,\infty)}
	h\cdot\mathbf 1_{[0,c]}\dd\lambda\right)
	= \int_{[0,\infty)} h\dd\lambda\tag{$\ast$}
    \end{equation*}
    To finish, $g_n$ is an increasing sequence
    (because $1 + x/n$ is decreasing in $n$), so by monotone convergence,
    \begin{equation*}
	\lim_{n\to\infty}\left(\int_0^\infty f_n\dd x\right)
	= \lim_{n\to\infty}\left(\int_0^\infty g_n\dd x\right)
	\overset{(\ast)}= \lim_{n\to\infty}\left(\int_{[0,\infty)} g_n\dd\lambda\right)
	\overset{\text{MCT}}= \int_{[0,\infty)} g\dd\lambda
	\overset{(\ast)}= \int_0^\infty g\dd x \overset{\text{(a)}}= \boxed\half.
    \end{equation*}

    (c) Suppose there exists such a function $F\colon[0,\infty)\to[0,\infty]$.
    Let $\widetilde F(x) = F$ at all $x\in X$ such that
    $F(x) \ge f_n(x)$ and $\widetilde F(x) = \infty$ for all other $x$.
    Since $F$ and $\widetilde F$ only disagree on
    a set of measure zero by hypothesis, they have the same (finite) integral,
    and $\widetilde F\ge f_n = \abs{f_n}$ for all $n$;
    thus by the dominated convergence theorem
    (where $\widetilde F$ dominates the sequence $f_n$),
    \begin{equation*}
	\half \overset{\text{(b)}}=
	\lim_{n\to\infty}\left(\int_{[0,\infty)} f_n\dd\lambda\right)
	= \int_{[0,\infty)} f\dd\lambda
	\overset{(\ast)}= \int_0^\infty f\dd x \overset{\text{(a)}}= 0.
    \end{equation*}
    This is a contradiction, so $F$ cannot exist.
\end{soln}

\pagebreak

\begin{problem}
    Show that the function
    \begin{equation*}
	f\colon[0,\infty)\to\RR,\quad
	f(x) = \begin{cases}
	    \frac{\sin x}x & x > 0\\
	    1 & x = 0
	\end{cases}
    \end{equation*}
    has an improper Riemann integral over $[0,\infty)$,
    but is not Lebesgue integrable on $[0,\infty)$.
\end{problem}

\begin{soln}
    Let $I_n = [(n-1)\pi, n\pi]$ and
    \begin{equation*}
	a_n = \int_{I_n} f\dd\lambda = \int_{(n-1)\pi}^{n\pi} f\dd x,
    \end{equation*}
    noting that $f$ is continuous hence Riemann integrable on each $I_n$.
    For $n > 1$, we can bound
    \begin{equation*}
	\tag{$\ast\ast$}
	\frac 1{n\pi} = \abs{\int_{(n-1)\pi}^{n\pi} \frac{\sin x}{n\pi}\dd x}
	\le \underbrace{\abs{\int_{(n-1)\pi}^{n\pi} \frac{\sin x}x\dd x}}_{\abs{a_n}}
	\le \abs{\int_{(n-1)\pi}^{n\pi} \frac{\sin x}{(n-1)\pi}\dd x}
	= \frac 1{(n-1)\pi}
    \end{equation*}
    with the lower bound holding for $n=1$ as well.
    Since $f$ is nonnegative exactly when $\sin x\ge0$, equivalently
    $x\in I_{2n-1}$ for some $n\in\NN$
    (and nonpositive if and only if $x\in I_{2n}$
    for some $n\in\NN$),
    \begin{equation*}
	f_+ = \sum_{n=1}^\infty f\cdot\mathbf 1_{I_{2n-1}},\quad
	f_- = \sum_{n=1}^\infty f\cdot\mathbf 1_{I_{2n}}.
    \end{equation*}
    In particular, by HW04 \#5a,
    \begin{equation*}
	\int_{[0,\infty)} f_+\dd\lambda
	= \sum_{n=0}^\infty a_{2n+1} \ge \sum_{n=0}^\infty \frac{1}{(2n+1)\pi}
	= \infty,\quad
	\int_{[0,\infty)} f_-\dd\lambda
	= \sum_{n=1}^\infty a_{2n} \ge \sum_{n=1}^\infty \frac{1}{2n\pi}
	= \infty.
    \end{equation*}
    Thus $\int f_+\dd\lambda$ and $\int f_-\dd\lambda$ are both $\infty$,
    so $f$ is not Lebesgue integrable on $[0,\infty)$.

    For Riemann integrability, define $F\colon[0,\infty)\to\RR$ by
    \begin{equation*}
	F(y) = \int_0^y f\dd x.
    \end{equation*}
    Given $c\in[0,\infty)$, let $n$ be the largest integer
    such that $c\in I_n$ (so if $k\in\NN$ such that $c=k\pi$ then $n = k+1$).
    Since $f$ is either always nonpositive
    or always nonnegative on $I_n$, $F$ is monotonic on $I_n$,
    meaning $F(c)$ lies between $F(n)$ and $F(n+1)$. Thus
    \begin{equation*}
	\lim_{c\to\infty}F(c) = \lim_{n\to\infty}F(n)
	= \sum_{k=1}^\infty a_k = \sum_{k=1}^\infty (-1)^{k+1}\abs{a_k}.
    \end{equation*}
    Since $\abs{a_k} \le \frac 1{(k-1)\pi} \le \abs{a_{k-1}}$ by $(\ast\ast)$
    for all $k\in\NN$, the sequence $\abs{a_k}$
    converges to $0$, so the sum on the right-hand side converges
    by the alternating series test. This means
    \begin{equation*}
	\lim_{c\to\infty}F(c) = \int_0^\infty f\dd x
    \end{equation*}
    is finite, hence $f$ has an improper Riemann integral on $[0,\infty)$.
\end{soln}		










\end{document}
